\scanpage{163}\mspage{21}
% archive sheet number: 163
\deleted{commutatif, d'où \deleted{\unsure{dém.}}\inserted{$\mathrm{PS}1$} \unsure{Voilà}. Pour voir que $b)\Rightarrow\mathrm{PS}1$ \unsure{vient} \unsure{au-dessous de b|au-dessus sur le} \deleted{\unsure{est localement} \illegible{2}} cf 2.5.}

\uline{Corollaire} 2.8. \deleted{Soit $\unsure{\mathscr X|\mathcal X}$ \illegible{4}, topos, $S\subset\mathscr X^{\wedge}$ \illegible{2} \unsure{petite}.} Soient $\mathscr X,S$ comme ds 2.6. \deleted{équivalentes, les condits. équivalentes} Cond. équiv :

a) $\mathscr X$ est loc. irréductible, i.e. les topos induits $\mathscr X_{/U}$, où $U$ est $0$-connexe, sont irréductibles

\deleted{\unsure{b)}} \unsure{b)} \deleted{\unsure{Si} \illegible{1}} \unsure{Tout} produit fibré $U\times_W V$ d'objets \deleted{\unsure{non}}\inserted{$0$-connexes} \deleted{$U,V$ non \unsure{vides}} ``non vides''.

\inserted{b') Tout produit fibré $U\times_W V$, avec $U,V$ ``non vides'', $W$ $0$-connexe, est ``non vide''.}

c) les $\mathscr X_{/a}$ avec $a\in S$ sont irréductibles

d) Comme b), avec $U,V,W$ dans $S$.

e) $S$ est \deleted{loc. fbt}\inserted{loc. ps.} filtrante à g \inserted{(= loc. fbt. filtrante à g.)}, i.e satisfait $\mathrm{PS}*1$.

Ici les implications $\pm$ évidentes sont

\diagram{\deleted{\begin{tikzcd}[ampersand replacement=\&,row sep=small,column sep=small]
\&b\arrow[r,Rightarrow]\arrow[d,Rightarrow]\&a\arrow[d,Rightarrow]\\
\unsure{e}\arrow[r,Leftrightarrow]\&d\arrow[r,Rightarrow]\arrow[dr,Leftrightarrow]\&\unsure{c}\arrow[d,Rightarrow]\\
\&\&\unsure{e}
\end{tikzcd}}}

\diagram{\begin{tikzcd}[ampersand replacement=\&,row sep=small,column sep=small]
b'\arrow[r,Rightarrow]\&b\arrow[rr,Rightarrow]\arrow[d,Rightarrow]\&\&a\arrow[d,Rightarrow]\\
\&d\arrow[r,Leftrightarrow]\&e\&c\arrow[l,Rightarrow]
\end{tikzcd}}

Prouvons les implications de la ligne du bas. \inserted{Considérons d),} \unsure{Si|Soit}
\diagram{\begin{tikzcd}[ampersand replacement=\&,row sep=small,column sep=small]
a\arrow[dr]\&\\
\&x\\
b\arrow[ur]\&
\end{tikzcd}}
dans $S$, dire que le produit fibré dans $\mathscr X$ est ``non vide'' signifie, puisque $S$ est un syst. de générateurs \deleted{\unsure{non vides}} et en particulier \uline{non vides}, qu'il existe $c$ dans $S$ et $c\longrightarrow a\times_x b$, ou ce qui revient au $=$,
\diagram{\begin{tikzcd}[ampersand replacement=\&,row sep=small,column sep=small]
\&a\\
c\arrow[ur,dashed]\arrow[dr,dashed]\&\\
\&b
\end{tikzcd}}
rendant commutatif le diagramme
\diagram{\begin{tikzcd}[ampersand replacement=\&,row sep=small,column sep=small]
\&a\arrow[dr]\&\\
c\arrow[ur,dashed]\arrow[dr,dashed]\&\&x.\\
\&b\arrow[ur]\&
\end{tikzcd}}
Donc d) équivaut à $\mathrm{PS}1$, i.e : e) ($\mathscr X$ loc. fbt. filtrante à g).

Considérons la condition c) : les $\mathscr X_{/a}=(S_{/a})^{\sim}$ sont irréductibles. Cela signifie que pour deux sous-objets \deleted{\unsure{ouverts}}\inserted{\uline{non vides}} $U,V\subset a$ dans $\mathscr X^{\wedge}$, on a $U\cap V\ne\varnothing$, ou encore que pour deux cribles
